% Transcription — fonds Alexandre Grothendieck, shelfmark 6, batch 4 (pages 61-66).
% Established from the facsimile archives/batches/6/batch-04.pdf, cut from a
% copy of 6.pdf fetched through the project's own relay
% (https://grothendieck-relay.onrender.com/source/6.pdf) rather than from
% grothendieck.umontpellier.fr directly; per the skill transcribe-grothendieck.
% The batch file carries no cover sheet: PDF page k is the archivists' page
% 60+k. This is the folder's last batch (six pages). The archivists' pencil
% numbers bottom-left do NOT agree with the facsimile here: pages 61-66 are
% pencilled 72, 73, 74, 75, 76, 77 (the same re-ordering after pencilling
% that batch 2 records from its p. 39, pencilled 50). \page{} follows the
% facsimile's order; the pencil numbers are recorded here only.
%
% Pass: Opus 5.5 (claude-opus-5-5), 2026-09-26 — first pass, unchecked against the pages by a human.
%
% The folder sits in the inventory group "4-9" « Cristaux »; it has its own
% date in the catalogue, which \dating{} carries verbatim, with no group
% sentence.
%
% What the pages are, and the authorship decisions.
% - Pages 61, 62, 64, 65: HIS OWN MANUSCRIPT NOTES (the folder's « notes
%   manuscrites (s.d.) »), transcribed in full. Evidence: the same fast
%   cursive, blue ink, as his ink corrections of the typed letter (batch 1)
%   and his manuscript notes of batch 2 (pp. 33-36, 40); his abbreviations
%   (« B.T. », « t.m. », « loc. », « t.f. », « isom. », « hom. », « NB »),
%   running self-corrections (struck words with interlinear replacements
%   throughout), and a marginal aside in his idiom, « Zar ou fppf, c'est
%   kif-kif » (p. 65). They are working notes, not a fair copy; no date on
%   them. Nothing on these pages is typed, and nothing is in another hand.
% - The folder's « tapuscrit (s.d.) » does not fall on these pages (it is
%   presumably the typescript material of batch 3); no page here is someone
%   else's document, so no page is reduced to a \note.
%
% Pages skipped (no \page): 63, a leaf with nothing of his — only the
% show-through of p. 64 (no ink of its own, checked at raised contrast);
% 66, a blank sheet (printed letterhead only, nothing written).
%
% Pages summarised: none.
%
% His pagination: none on these leaves.
%
% What the batch is about. Continuation of notes (begun before p. 61, in
% batch 3) on the Dieudonné crystal D^*(G) of truncated Barsotti-Tate groups
% and its dual description through the « triangle exact de dualité »
% relating the complexes l^G, l^{G*}[1] and Delta(G). p. 61: for G étale,
% l^G = 0 and the triangle gives Delta(G) ≃ l^{G*}[1] ≃ epsilon^*(M/p) ≃
% RHom_Z(G, O_S) ≃ (G^v ⊗^L F_p) ⊗ O_S, so no extra structure on the crystal;
% groups of multiplicative type (killed by a power of p), defined by the
% étale G^*: D^*(G) = Theta_*(G^*), V an isomorphism, F then determined by
% V; equivalence « groupes de B.T. tronqués d'échelon n de t.m. ≃ cristaux
% loc. libres de t.f. sur S/Lambda_n avec V isom. (et F nilpotent) », and
% analogously for multiplicative type (p. 62), where now l^{G*}[1] = 0 and
% l^G ≃ G^* ⊗^L O_S. Homomorphisms between étale p-groups and groups of
% multiplicative type (a: nul, via F^n for n large; b: left as a heading).
% pp. 64-65: over S of char. p, the classical equivalence between finite
% étale group schemes killed by p (≃ locally free F_p-sheaves) and locally
% free O_S-modules M with F_M : M^(p) ⥲ M (unit-root F-modules): full
% faithfulness by descent to the Artin-Schreier sequence 0 -> Z/pZ -> G_a
% -> G_a; essential surjectivity via the kernel of psi = F_M - id, reduced
% to the Lemma over an algebraically (corrected to « sép. ») closed field:
% E ⊗_{F_p} k ⥲ M. Then the start of the Witt-vector version: C finite
% étale p-groups, C' sheaves of W_S-modules locally sums of W_{n,S} with
% F_M, C'' smooth W_S-module schemes, the functors phi, i, psi = E ⊗ W_S;
% the page breaks off (« représente le faisceau... »).
%
% Notation fixed once for the batch: his « l. » and « Delta. » (index a
% lower dot) as \ell_{\cdot}, \Delta_{\cdot}; the dual with a check as
% \check{\ell}, \check{G}; his hollow D as \mathbb{D}^*; RHom as
% \mathbb{R}\underline{\mathrm{Hom}}; derived tensor as \overset{\mathbb{L}}{\otimes};
% \underline{O}_S as in batches 1-2; the curly C of the categories as
% \mathcal{C}; the Witt sheaves as W_S, W_{n,S}, with his superscripts
% « zar » and (read with doubt) « fl ».
%
% Apparatus: 9 \ill{} (5 of them inside \struck{}), 9 \uncertain{}.
%
% Readings to check first: « en sus » and the scribble after it, p. 61;
% « prolongeant » (under-bracket, p. 61); « Lambda_n » and « nilpotent »,
% foot of p. 61; the mark above « tout hom. » and « d'où \ill », p. 62;
% « \ill{} de F_p-vectoriels » (after « faisceaux »), p. 64; the label of
% psi and « épimorphisme », p. 64; « en fait », p. 64; the superscript
% « fl » of W_S in the C'' line, p. 65.

\documentclass[11pt,a4paper]{article}
% Le préambule commun à toutes les transcriptions du fonds Grothendieck.
%
% Il tient en une idée : **l'incertitude se compose, elle ne se résout pas.**
% Une transcription qui lisse un mot douteux a détruit la seule chose pour
% laquelle on la faisait — un lecteur doit pouvoir savoir, à chaque endroit,
% ce qui était sur le feuillet et ce qui a été ajouté. Les six macros qui
% suivent sont donc l'appareil critique, et non de la décoration.
%
% Les mêmes commandes sont reconnues par `scripts/render.mjs`, qui produit la
% vue de lecture du site. Une macro ajoutée ici doit l'être aussi là-bas,
% faute de quoi le rendu échouera bruyamment — ce qui est voulu : un rendu
% silencieusement faux est pire qu'un rendu absent.

\ProvidesPackage{grothendieck}[2026/08/08 Transcriptions du fonds Grothendieck]

\usepackage{amsmath,amssymb,amsthm}
\usepackage{mathrsfs}
\usepackage{stmaryrd}
% Les diagrammes commutatifs sont partout dans ce fonds : c'est la langue
% même de la géométrie algébrique de Grothendieck, et une transcription qui
% les rendrait en prose ne transcrirait rien.
\usepackage{tikz-cd}
% Français par défaut — c'est la langue des feuillets — mais l'anglais est
% chargé aussi : la traduction et le résumé sont des documents anglais, et la
% typographie française (espaces avant les deux-points) y serait une faute.
% Ces documents basculent par \selectlanguage{english} en tête de corps.
\usepackage[english,french]{babel}
\usepackage{fontspec}
\usepackage{geometry}
\geometry{a4paper,margin=25mm,marginparwidth=28mm}
\usepackage{marginnote}
\usepackage{xcolor}
% ulem, not soul. `soul`'s \st and \ul are built on a character-by-character
% scan that XeLaTeX's font machinery breaks: the document fails with an
% undefined control sequence rather than degrading. ulem does the same two jobs
% and is Unicode-engine safe. `normalem` stops it hijacking \emph, which the
% transcriptions use for Grothendieck's own emphasis.
\usepackage[normalem]{ulem}
\usepackage{hyperref}
\hypersetup{colorlinks=true,linkcolor=black,urlcolor=black,pdfborder={0 0 0}}

% --- Métadonnées du lot -----------------------------------------------------
% Reprises telles quelles par le script de rendu, qui les affiche en tête de
% la vue de lecture. Elles fixent ce que le fichier prétend couvrir : sans
% elles, une transcription flotte sans point d'attache dans un fonds de seize
% mille feuillets.

\newcommand{\folder}[1]{\gdef\grothFolder{#1}}
\newcommand{\batch}[1]{\gdef\grothBatch{#1}}
\newcommand{\leaves}[2]{\gdef\grothFirst{#1}\gdef\grothLast{#2}}
\newcommand{\foldertitle}[1]{\gdef\grothFoldertitle{#1}}
\gdef\grothFolder{?}\gdef\grothBatch{?}\gdef\grothFirst{?}\gdef\grothLast{?}\gdef\grothFoldertitle{}

% --- Le résumé --------------------------------------------------------------
%
% Il ouvre la lecture modernisée, sous le titre « Résumé », et il est composé
% en petit. Ce n'est pas un ornement : c'est le seul passage du document qui
% ne soit pas des mathématiques, et un lecteur doit pouvoir le reconnaître
% avant de l'avoir lu — puis le sauter s'il connaît déjà le sujet, ou s'y
% arrêter s'il ne le connaît pas. Le filet qui le suit marque la reprise du
% corps.

\newenvironment{resume}
  {\small\setlength{\parskip}{0.45em}}
  {\par\medskip\hrule\medskip}

% --- Mots-clés ---------------------------------------------------------------
%
% En anglais, en clôture du résumé de la lecture modernisée : le vocabulaire
% moderne sous lequel chercher ce que les feuillets construisent. Le manifeste
% du site (`npm run manifest`) les extrait pour étiqueter la cote entière —
% cette ligne est la source unique des tags, il n'y a pas de fichier à côté.
\newcommand{\keywords}[1]{\par\smallskip\noindent
  {\footnotesize\textsc{Keywords} --- \textit{#1}}\par}

% --- L'appareil critique ----------------------------------------------------

% Le numéro de feuillet, en marge. C'est l'ancre qui permet de régler un
% désaccord sur une formule en regardant *un* feuillet plutôt qu'une chemise
% de six cents. Le site s'en sert aussi pour tourner les pages du fac-similé
% quand on fait défiler la transcription.
\newcommand{\leaf}[1]{%
  \marginnote{\footnotesize\color{gray!70}\textsf{#1}}%
  \label{leaf:#1}\ignorespaces}

% Illisible : on ne devine pas. Un mot inventé qui se lit comme les autres est
% le pire résultat possible.
\newcommand{\ill}{\textcolor{red!60!black}{[\,\ldots\,]}}

% Lecture douteuse : le mot est donné, et signalé comme incertain.
\newcommand{\uncertain}[1]{\uline{#1}}

% Ajout de l'éditeur — une lettre manquante, un mot sous-entendu.
\newcommand{\add}[1]{\textcolor{blue!50!black}{[#1]}}

% Biffé par Grothendieck lui-même. Ce qu'il a rayé fait partie du document :
% une rature dit souvent où la pensée a changé de direction.
\newcommand{\struck}[1]{\sout{#1}}

% Note du transcripteur, sur l'état matériel du feuillet.
\newcommand{\note}[1]{\par\smallskip{\small\color{gray!60!black}\textit{[#1]}}\par\smallskip}

% Note marginale de Grothendieck, distincte du corps du texte.
\newcommand{\marginal}[1]{\par\smallskip\noindent\hspace{1em}%
  {\small\color{gray!60!black}#1}\par\smallskip}

% --- Titre ------------------------------------------------------------------

\AtBeginDocument{%
  \begin{center}
    {\large\bfseries Fonds Alexandre Grothendieck --- cote n\textsuperscript{o}~\grothFolder}\\[2pt]
    {\itshape\grothFoldertitle}\\[4pt]
    {\small feuillets \grothFirst--\grothLast\ (lot \grothBatch)}\\[2pt]
    {\footnotesize\color{gray!60!black}Transcription établie d'après le fac-similé numérisé
     par l'Université de Montpellier.\\
     Les lectures incertaines sont soulignées, les illisibles notées [\,\ldots\,],
     les ajouts entre crochets.}
  \end{center}
  \vspace{1em}}

\endinput


\folder{6}
\batch{4}
\pages{61}{66}
\dating{1966-[à partir de 1970]}
\watermark{Édition de démonstration}
\foldertitle{Lettre Tate (mai 66) (Cristaux) : lettre (1966), tapuscrit (s.d.), notes manuscrites (s.d.)}

\begin{document}

\page{61}
\note{notes manuscrites, à l'encre, de sa main ; la page commence au milieu
d'un argument commencé plus haut dans le dossier (lot précédent).}
Comme $\ell^{G}_{\cdot}=0$, le triangle exact de dualité pour $G$ se réduit
à :

\begin{tikzcd}
 & \check{\ell}^{G^*}_{\cdot}[1] \arrow[dl] & \\
0 \arrow[rr] & & \Delta_{\cdot}(G) \arrow[ul, "\wr"']
\end{tikzcd}

\note{la flèche montante porte un petit signe, lu comme $\wr$
(isomorphisme).}
i.e.\ ne donne pas de structure supplémentaire \struck{\ill{}} en sus
\ill{} du cristal $\mathcal{M}$.\note{« en sus » est écrit au-dessus du mot
biffé ; le signe qui suit est surchargé.} On a
\[
  \check{\ell}^{G^*}_{\cdot}[1]\simeq\Delta^{*}_{\cdot}(G)\simeq
  \varepsilon^{*}(\mathcal{M}_{/p})=
  \mathbb{R}\underline{\mathrm{Hom}}_{\mathbb{Z}}(G,\underline{O}_S)
  \simeq\check{G}\overset{\mathbb{L}}{\otimes}_{\mathbb{Z}}\underline{O}_S
  \simeq(\check{G}\overset{\mathbb{L}}{\otimes}_{\mathbb{Z}}\mathbb{F}_p)
  \otimes_{\mathbb{F}_p}\underline{O}_S\ .
\]
\note{l'astérisque de $\Delta^{*}_{\cdot}$ est surchargé sur le $\Delta$ ;
devant $\mathbb{R}\underline{\mathrm{Hom}}$, un signe noirci, biffé.}

\emph{Groupes de type multiplicatif} (annulés par puiss.\ de $p$)

Soit $G$ \struck{\ill{}} un tel, il est défini par $G^*$, groupe étale annulé
par $p^n$, auquel s'appliquent les réflexions précédentes. On a
\uncertain{donc}
\[
  \mathbb{D}^*(G)=\Theta_{*}(G^*)
\]
i.e.
\[
  \mathbb{D}^*(G)_{S'}=(G^*_{S'})\otimes_{\mathbb{Z}_p}\underline{O}_{S'}
\]
\note{sous $G^*_{S'}$, une accolade renvoie à : « groupe étale sur $S'$
\uncertain{prolongeant} $G^*$ ».}

Ici
\[
  V\colon\mathbb{D}^*(G)\longrightarrow\mathbb{D}^*(G^{(p)})=\mathbb{D}^*(G)^{(p)}
\]
est un isomorphisme, savoir
$\Theta_{*}(F_{G^*}\colon G^*\to G^{*(p)})$, et $F$ est par suite déterminé
uniquement en termes de $V$.\note{« par suite » est ajouté au-dessus de la
ligne.}

On a encore des équivalences de catégories
\begin{quote}
groupes de B.T.\ tronqués d'échelon $n$ et de t.m.\ $\;\approx\;$ cristaux
loc.\ libres de t.f.\ sur $S/$\uncertain{$\Lambda_n$}, avec $V$ un isom
(et $F$ \uncertain{nilpotent})
\end{quote}

\page{62}
et équivalences analogues pour les groupes \struck{de B.T.\ tr} de type
multiplicatif annulés par une puissance de $p$. \struck{On conclut encore}
Ici $\check{\ell}^{G^*}_{\cdot}[1]=0$, et le triangle exact de dualité se
réduit à :

\begin{tikzcd}
 & 0 \arrow[dl] & \\
\ell^{G}_{\cdot} \arrow[rr] & & \Delta_{\cdot}(G) \arrow[ul]
\end{tikzcd}

avec $\ell^{G}_{\cdot}$ donné comme dessus
\[
  \ell^{G}_{\cdot}\simeq G^*\overset{\mathbb{L}}{\otimes}_{\mathbb{Z}}\underline{O}_S
  \simeq(G^*\overset{\mathbb{L}}{\otimes}_{\mathbb{Z}}\mathbb{F}_p)
  \otimes_{\mathbb{F}_p}\underline{O}_S\ .
\]

\emph{Homomorphismes entre $p$-groupes étales et groupes de t.m.}
\struck{\uncertain{connexes}}

a) Un hom.\ d'un groupe $G$ de t.m.\ dans un étale $G'$ est nul. D'autre
part, tout \ill{} hom.\ de $F$-$V$-cristaux (ou simplement de $F$-cristaux,
ou de $V$-cristaux) de $M'$ dans $M$ est nul.\note{les lettres $G$ et $G'$
sont écrites au-dessus de « groupe » et de « étale » ; au-dessus de « tout
hom. », un petit signe non lu.} En effet, écrivons la
\uncertain{commutation} à $F^n$ pour $n$ assez grand, en notant que
$F^n_{M'}$ est un isom et $F^n_{M}$ est nul (pour $n$ grand), d'où \ill{} le
résultat pour les $F$-hom., et on procède de même pour les $V$-hom.

b) Hom.\ d'un groupe étale $G$ dans un groupe de t.m.\ $G'$.

\page{64}
$S$ schéma de car $p$.

$\mathcal{C}$ catégorie des schémas en groupes finis étales $E$ sur $S$
annulés par $p$

$\approx$ catégorie des faisceaux \ill{} de $\mathbb{F}_p$-vectoriels
loc.\ libres de type fini\note{le nom $E$ est ajouté au-dessus de la ligne,
après « étales ».}

$\mathcal{C}'$ catégorie des Modules loc.\ libres $M$ sur $S$, munis de
\[
  F_M\colon M^{(p)}\xrightarrow{\ \sim\ }M
\]

$\mathcal{C}\xrightarrow{\ \varphi\ }\mathcal{C}'$
\[
  E\longmapsto(E\otimes_{\mathbb{F}_p}\underline{O}_S\,,\;
  F_E\otimes_{\mathbb{F}_p}\underline{O}_S)
\]
\emph{NB}
$\bigl((E\otimes_{\mathbb{F}_p}\underline{O}_S)^{(p)}\simeq
E^{(p)}\otimes_{\mathbb{F}_p}\underline{O}_S
\xleftarrow[\ F_{E/S}\ ]{\ \sim\ }E\otimes_{\mathbb{F}_p}\underline{O}_S\bigr)$

\emph{Proposition} \emph{Le foncteur $\varphi$ est une équivalence de
catégories.}

a) \emph{Pleinement fidèle.} Par passage \uncertain{aux}
$\underline{\mathrm{Hom}}$, revient à ceci :
\[
  \Gamma(S,E)\xrightarrow{\ \sim\ }\bigl\lbrace x\in\Gamma(S,E\otimes_{\mathbb{F}_p}
  \underline{O}_S)\ \big|\ F(x^{(p)})=x\bigr\rbrace
\]
Par descente, on se ramène au cas $E$ constant, puis
$E=\mathbb{F}_{p\,S}$, et on est ramené alors à l'assertion que
\[
  0\to(\mathbb{Z}/p\mathbb{Z})_S\to\mathbb{G}_{a,S}
  \xrightarrow{\ x\mapsto x^p-x\ }\mathbb{G}_{a,S}
\]
est exact, ce qui est bien connu (il suffit de le voir sur $\mathbb{F}_p$
par changement de base \ldots.)

b) \emph{Essentiellement surjectif.} Si $(M,F)$ est donné, il suffit de voir
qu'il est \emph{loc.} de la forme voulue. Considérons
\[
  M\xrightarrow{\ \psi\ }M\ ,\qquad \psi x=F_M\,x-x\ ;
\]
\note{l'étiquette $\psi x=F_M\,x-x$ est au-dessus de la flèche ; entre
$F_M$ et $x$ un signe resserré, peut-être l'exposant $(p)$ de $x^{(p)}$.}
on voit que c'est un \uncertain{épimorphisme} étale de schémas en groupes,
donc son noyau est un schéma en groupes étale $E$ sur $S$.\note{le début de
« épimorphisme » est surchargé.} Pour montrer qu'il est fini sur $S$, il
suffit de prouver qu'il est de $\mathbb{F}_p$-rang \struck{ég} localement
constant --- \uncertain{en fait}, égal au rang de $M$.\note{« $\mathbb{F}_p$- »
est ajouté au-dessus de « rang ».} Pour prouver ensuite que
$E\otimes_{\mathbb{F}_p}\underline{O}_S\to M$ est un isom, il suffit de le
prouver fibre par fibre (géom.). Bref, on est ramené au cas

\page{65}
d'un corps de base alg.\ clos, et à prouver ceci :

\emph{Lemme} Soit $M$ vectoriel sur $k$ \struck{\ill{}} sép.\ clos de car
$p$, muni de $F_M\colon M^{(p)}\xrightarrow{\sim}M$. Soit
$E=\lbrace x\in M,\ F x^{(p)}=x\rbrace$. Alors
$E\otimes_{\mathbb{F}_p}k\xrightarrow{\sim}M$.\note{« sép. » est écrit
au-dessus du mot biffé ; l'énoncé est marqué d'un trait vertical dans la
marge ; il est suivi d'un double trait de séparation.}

$\mathcal{C}$ catégorie des schémas en $p$-groupes \struck{finis loc.\
libres} finis étales sur $S$

i.e.\ des faisceaux en $p$-groupes loc.\ constants de t.f.

\struck{\ill{}}\note{« $p$- », « finis étales » et « constants de t.f. »
sont ajoutés au-dessus ou au-dessous de la ligne ; la ligne suivante est
entièrement biffée, illisible.}

\marginal{Zar ou fppf, c'est kif-kif}\note{dans la marge de gauche, avec
une accolade, en regard de la définition de $\mathcal{C}'$.}
$\mathcal{C}'$ catégorie des \struck{schémas} faisceaux en
$W_S^{\mathrm{zar}}$-modules $M$ loc.\ isom.\ à une somme de copies de
$W_{n,S}$, munis de \struck{\ill{}}
$M^{(p)}\xrightarrow[\ \sim\ ]{\ F_M\ }M$

$\mathcal{C}''$ catégorie des schémas en
$W_S^{\mathrm{fl}}$-modules lisses\note{l'exposant « fl » de $W_S$ est lu
avec doute.}

\begin{tikzcd}
\mathcal{C} \arrow[r, "\varphi"] \arrow[dr, "\psi"'] & \mathcal{C}' \arrow[d, "i\ \sim"] \\
 & \mathcal{C}''
\end{tikzcd}

$\varphi(E)=E\otimes_{\mathbb{Z}_p}W_S^{\mathrm{zar}}$ avec le $F$ évident

$i(M)=$ \struck{faisceau repr} schéma qui représente le faisceau\ldots

$\psi(E)=E\otimes_{\mathbb{Z}_p}W_S$\note{« schéma qui » est écrit
au-dessus des mots biffés. La note s'arrête là ; le reste de la page est
vide.}

\end{document}
